\section{Reproducibility}

The rescue package binds results to a canonical JSON evidence registry. Each source report is identified by a project-relative path and SHA-256. The registry records dataset version, partition, evidence class, metric contract, seed count, values, confidence intervals, and limitations. Paper tables are generated from that registry rather than transcribed manually.

The repository includes the Loop191, Loop192, and Loop206 configs; model constructors; preprocessing and postprocessing implementations; metric code; Loop205 regional saliency protocol; Loop206 active-contour implementation; evidence builder; table builder; tests; and website source. Runtime data, checkpoints, probability caches, and large artifacts remain outside Git. The selected Loop206 control and candidate checkpoints are verified by SHA-256 at application startup.

For arbitrary uploaded images, the candidate model requires a deployment saliency prior fitted only from the 308 Loop206 fit groups. Deployment is authorized only if the resulting contour generator reproduces all 76 frozen holdout clean contours byte-for-byte. If data identity, source hashes, model hashes, or contour parity fail, candidate upload inference remains disabled rather than using an approximate channel.

The recorded hardware is an RTX 5060 Ti workstation with Python 3.12. Environment dependencies are locked with \texttt{uv}. A second Windows laptop performs clean-clone CPU tests and LaTeX compilation without datasets or weights. The final release records test, paper-audit, PDF-build, GPU-smoke, and public-tunnel receipts.
